\documentclass[]{iat}
%Diese Beiden Pakete werden nur für das Beispieldokument benötigt und sollten von Ihnen gelöscht werden
\usepackage{listings}
\usepackage{scrhack}
%Hier können weitere benötigte Pakete Eingebunden werden

%Namen eingeben / Insert Name
\renewcommand{\author}{Li Hongfeng}
%Art der Arbeit / Scope (Project, Thesis...)
\providecommand{\scope}{Projektarbeit}
%Thema der Arbeit / Theme of Thesis
\renewcommand{\subject}{Application of Data Analytics in Parameter Monitoring and Optimization of WWTP}
%Schlagwörter / Keywords
\providecommand{\keywords}{Key_1, Key_2}
%Literaturliste *.bib / Bibliography
\providecommand{\bibfile}{IAT_Beispiel_literature}
%Matrikelnummer / Student ID
\providecommand{\studentID}{6104950}
%Betreuer /& Tutors
\providecommand{\tutora}{M.Sc. Yuying Shao}
%Prüfer / Examiner
\providecommand{\examinera}{Prof.Dr.-Ing. Kai Michels}


\hypersetup{%
	pdftitle	={\subject -- \author -- \today},
	pdfauthor	={\author},
	pdfsubject	={\subject},
	pdfkeywords ={\keywords}
}

\bibliography{\bibfile} 

\setlength{\footheight}{21pt}

\begin{document}
	\lstset{literate={ä}{{\"a}}1}
%Pfad zu Grafiken:
	\graphicspath{{./project_graphics/}}
% Sprachauswahl /Language Selection (ngerman/english)
	\selectlanguage{ngerman}
\pagenumbering{roman}
\input{iat_title}
%Urherberrechtserklärung / Confirmation of Conformity Comment if not needed, Select correct Language
\input{declaration_conformity}
\pagestyle{iat}
\tableofcontents
\chapter{Einleitung / Introduction}
\pagenumbering{arabic}
\setcounter{page}{1}
Data play a key role in contemporary research and can be collected to understand patterns and trends in the field under study. Specifically, data can be represented in the form of numbers, images, etc., to provide a comprehensive understanding of the process of inquiry, optimization of the process, and changes in trends.\\
\\
Data analysis and processing are fundamental to unlocking the potential of data. Through rigorous statistical analysis, we can identify important correlations and causal relationships that are not readily apparent.And the processing of the data can make the data analysis more clear and intuitive.\\
\\
Therefore data is necessary as an input for analytics, which is why the two are closely intertwined.
\section{Background}
Operational data of the WWTP are collected in real time through online monitoring systems. These data are preprocessed to remove obvious outliers and linearly interpolated to compensate for incidental flaws in the data collection process. The data that will be used in this paper is redundant and complex, so data analysis is particularly important, and the processed dataset is analyzed here using Python.\\
\\
The wastewater treatment process typically involves several key stages$^{[1]}$:\\
\\
The first step is pre-treatment, where large debris and gravel are removed in the pre-treatment step. Then primary treatment is carried out, using sedimentation method used to remove suspended solids from the wastewater, where the parameter "SS" plays an important role and the filtration effect depends on the variation of the concentration of suspended solids particles in the effluent.\\
\\
The next step is secondary treatment, which is most commonly used and is also known as activated sludge, where sludge and a pool of bacteria-rich sediment are added to the incoming wastewater and the mixture is agitated for several hours with an adequate supply of air. Suspended solids and many organic solids are absorbed or adsorbed by the sludge, while the organic matter is nitrified and denitrified by microorganisms. The amounts of air and sludge can be varied to control the level of treatment obtained. The sludge is then separated into a settling tank. The active biomass ("XBH" and "XBA") contained therein can be used to degrade organic matter and remove nutrients, which can be broken down by aerobic bacteria in the wastewater treatment system to remove wastes from the wastewater.\\
\\
Finally, there may be a tertiary treatment, a step aimed at filtering out nutrients and waste particles that can damage sensitive ecosystems; the wastewater passes through additional filter cells or ponds to remove more nutrients.
\section{Research Scope and Objectives}
This study includes the analysis of the relationships between various operational parameters in a wastewater treatment plant, such as determining the interactions and identifying significant dependencies between different parameters such as flow rate (FR), ammonia-nitrogen concentration (SNH), and dissolved oxygen concentration (SO), as well as the discovery of cyclic cycles in the treatment process that may correspond to daily, weekly, or seasonal variations. Repeated patterns or Motifs were also found in the operational data that could indicate regular operational cycles or recurring problems.
\section{Outline of the Thesis}
In this paper, we begin by outlining our methodology, detailing the specific plots and the 'stumpy' library used for data analysis. Then the data source is described, providing context and relevance to our study. Finally, we interpret the visual data, drawing conclusions from the analysis to highlight key insights and implications.
\chapter{Methodology}
\section{Significance of Data analysis}
Data analysis plays a critical role in the operation and optimization of wastewater treatment plants. Systematic analysis of historical operational data provides insight into the performance of the treatment process, identifying patterns and trends. Charts and visualizations are essential tools in data analysis. They provide a clear and intuitive way to present complex data, making it easier to understand and interpret information. Various graphs were generated in our study to visualize key parameters over time and identify important themes. The significance of these graphs will be discussed specifically in the following section.
\section{Tools and Libraries Used}
\subsection{Pandas$^{[2]}$}
Pandas is a widely-used library for data manipulation and analysis in Python. It provides data structures like DataFrames, which are convenient for handling structured data.\\
\\
First, it can read CSV data and load time-series data from a CSV file into a DataFrame. Second, it can perform operations such as setting up time indexes and slicing and dicing data for analysis. Then there is the ability to deal with missing data, and although it is not explicitly shown here, Pandas provides powerful ways to deal with missing or inconsistent data.
\subsection{STUMPY$^{[3]}$}
“STUMPY” is a powerful library for analyzing time series data, especially for matrix profile calculation. A matrix profile is a data structure used to label time series to show similarities between different subsequences. Motif discovery is mainly used to identify repeating patterns in time series data. By using the previously mentioned “stumpy” library, motifs can be detected and potential repeating sequences can be revealed to gain insights into the behavior and characteristics of the dataset. The significance of this technique is that it helps to identify important trends and anomalies, leading to a deeper understanding of the structure and meaning of the data. \\
\\
In this study case, STUMPY is used to compute matrix profiles for signals with different window sizes ("m1" and "m2"). And this window size serves to calculate the length of the time window for matrix summaries or other statistical properties. It is very important when analyzing time series because it determines the length of the subsequences to be compared. In the following analysis by performing Fourier transform or other spectral analysis of the time series to identify the major frequency components, and then converting the major frequency components to periods, then the two major periods are chosen as the window sizes m1 and m2.\\
\\
Among these matrix profiles can be helpful in identifying themes, i.e., the most similar subsequences in a time series. This type of analysis is essential to understand repeating patterns and anomalies in the data and can be used for anomaly detection, prediction etc. Theme discovery is mainly used to identify repeating patterns in time series data.
\subsection{Boxplot$^{[4]}$}
Boxplot is a statistical graph that can be used in this experiment to show the distribution and key statistical features of the WWTP data. It provides a concise snapshot of the data distribution by showing the minimum, first quartile (Q1), median (Q2), third quartile (Q3), and maximum values of the data, as well as possible outliers.\\
\\
In the data analysis below, the use of box-and-line plots helps:
A box-and-line plot allows you to visualize the extent of the data distribution and the centralized trend. Where the length of the box indicates the degree of dispersion of the data, while the position of the median line shows the central tendency of the data.\\
Boxplots clearly show outliers in the data, which are usually located at separate points on the outside of the box.\\
In addition boxplots can help spot symmetry and bias in the data. If the difference in whisker length between the two sides of the box is large, or the median is far from the center of the box, it means that the data may be skewed, which side-steps the anomalies in the data for this parameter in the wastewater treatment plant.
\subsection{FFT and Laplace transforms}
The Fourier Transform is used in data analysis to extract the frequency components of a signal and the period is calculated by the Laplace Transform. These two transforms have important applications in signal processing and data analysis.\\
\\
The main purpose of the Fourier Transform is to identify the periodic behavior in the signal. By finding the major frequencies in the spectrum, it is possible to understand the oscillatory properties and periodic characteristics of the system. Then, using the "output\_period" (output/period) function, the corresponding period can be calculated using the fact that in the Fourier function period and frequency are inverses of each other. The calculation formula is as follows:\\
\begin{equation}
period = 1/frequency
\end{equation}
Since here the data is sampled at 15-minute intervals, i.e., 4 times per hour, to convert the frequency to period, it is necessary to multiply by another 1/4 to reflect the hourly frequency.The specific formula is:\\
\begin{equation}
period = 1/frequency * 15/60 = 15/1440
\end{equation}
The Fourier transform helps to extract the frequency components of the signal from the time domain, thus identifying the major periodic behaviors. And the calculation of cycles through Laplace transform enables to quantify these periodic behaviors.
\subsection{Matplotlib$^{[5]}$}
Matplotlib is a plotting library in Python used for creating static, interactive, and animated visualizations.\\
\\
Matplotlib is used to plot time series data and can visualize time series data to identify patterns and trends. It is also possible to draw matrix profiles in plot aka visualize matrix profiles, view similarity scores of different subsequences as well as add annotations such as vertical lines to highlight.
\chapter{Comprehensive Overview of Wastewater Treatment}
\section{Data source of Wastewater Treatment Parameters}
In wastewater treatment plants, various parameters need to be monitored to ensure the efficiency and effectiveness of wastewater treatment. These parameters provide important information about the composition and quality of the wastewater and influence the design and operation of the treatment process. Understanding these parameters is critical to optimize the treatment process and ensure the safety of the wastewater discharged into the environment. The following sections detail the key parameters commonly measured at wastewater treatment plants, emphasizing their significance and role in the treatment workflow.\\
\\
In wastewater treatment plants, parameters such as Sl, SS, XI, XS, XBH, XBA, XP, SO, SNO, SNH, SND, XND, SALK, TSS, FR, T provide valuable information for understanding the treatment process. In this study, we used data analysis techniques to identify patterns and trends in these parameters for better management and optimization of wastewater treatment processes.\\
\\
SI is Soluble Inorganic, which indicates the soluble inorganic substances in the effluent, such as nitrates and chlorides. It primarily assesses the amount of inorganic substances in wastewater, which affects the efficiency of chemical and biological reactions in the wastewater treatment process.\\
\\
SS stands for Suspended Solids and indicates the concentration of suspended solid particles in the effluent, usually expressed in milligrams per liter (mg/L).The indicator is one of the most basic data of sewage, for the determination of suspended solids in the water and plant water, can be used to reflect the sewage through the primary sedimentation tank. Sedimentation tank treatment, the reduction of suspended solids, which is the main basis for reflecting the structure of the sedimentation efficiency and assess the effectiveness of the sedimentation tank and filtration system is an important indicator.\\
\\
XI is inert particulate organic matter, which is organic particulate matter in effluent that is not readily degradable. It can be used to assess the impact of difficult-to-degrade organic matter in effluent on the treatment system.\\
\\
XS can be understood as Slowly Biodegradable Substrates, i.e., organic matter in wastewater that is slowly degradable by microorganisms. It affects the rate of microbial degradation and treatment efficiency during biological treatment.\\
\\
XBH is the abbreviation of Active biomass, Heterotrophic, which generally refers to the heterotrophic microbial groups that feed on organic matter in sewage. These microorganisms play a vital role in the biological decomposition of organic matter and the reduction of BOD (Biochemical Oxygen Demand) and COD (Chemical Oxygen Demand) in wastewater treatment processes.\\
\\
XBA is the abbreviation of Active Biomass, Autotrophic, which refers to the group of autotrophic microorganisms in wastewater that use inorganic substances (e.g. ammonia, nitrate) as a nutrient source. Autotrophic bacteria oxidize ammonia to nitrate in the nitrification process, which is the key to the biological nitrogen removal process.This step is critical for removing nitrogen from wastewater and preventing eutrophication of receiving waters.\\
\\
SO stands for dissolved oxygen, which is the concentration of oxygen dissolved in sewage. It is a good reflection of the oxygen supply to the wastewater treatment system and is an important parameter of the aerobic treatment process.\\
\\
XP refers to Particulate Organic Matter, which is organic material in the form of particles in wastewater, which can affect the solid-liquid separation effect in the wastewater treatment process.\\
\\
SNO is Nitrate and Nitrite Nitrogen, specifically in effluent as nitrogen in the form of nitrates and nitrites. Its content allows monitoring the efficiency ofthe nitrification and denitrification reactions in the wastewater treatment process.\\
\\
SNH stands for Ammonium Nitrogen, which is found in effluent in the form of ammonia irons. Measurement of SNH allows for an assessment of the degree of ammonia contamination in the effluent and is the target substance for ammonia oxidation in the biological nitrogen removal process.\\
\\
SND is known as Soluble Organic Nitrogen, nitrogen in the form of dissolved organic matter in wastewater. It is used to assess the organic nitrogen content of wastewater and affects the efficiency of nitrogen removal during wastewater treatment.\\
\\
XND is known as Particulate Organic Nitrogen, which occurs as particulate organic nitrogen in effluent. It is mainly used to assess the effect of particulate organic nitrogen in effluent on the treatment process.\\
\\
SALK represents alkalinity, which is the concentration of carbonates, bicarbonates and other alkaline substances in the wastewater. It can affect the stabilization of pH in the wastewater treatment process and is an important parameter in the biological nitrogen removal process.\\
\\
The full name of TSS is Total Suspended Solids, which is expressed as the total concentration of suspended solids, both organic and inorganic, in wastewater and is defined as the dry weight of all suspended particles that are not in a dissolved state.The removal of TSS is usually achieved through physical treatment processes such as sedimentation, filtration, and flotation. It assesses the overall amount of solids in the effluent and influences solid-liquid separation and sedimentation.\\
\\
FR translates to Flow Rate, which is the amount of water that passes through the wastewater treatment system per unit of time. It is used to reflect the capacity and efficiency of the wastewater treatment system.\\
\\
T simply represents the Temperature of the wastewater, which also reflects the reaction rate of the wastewater treatment process.
\chapter{Result}
\section{Data time and index position settings}
The time range of data recording is set here as day 245 to day 609. Specify the time period of the data to be analyzed by defining the start time and end time as 245 and 609 days. Using the generated time arrays, time points were added to the data frame and indexed for time series analysis. In this way, each data point corresponds to a specific time point, ensuring the accuracy of the analysis process. This timeframe was chosen to capture the performance of the WWTP specifically over time to provide a comprehensive basis for analysis.\\
\\
The process of determining the index positions of the smallest values in the matrices 'mp1' and 'mp2' in the following can be converted to time points in days by dividing these index positions by 4 * 24, i.e., 4 data points per hour, 24 hours per day. The identification of these patterns helps to understand the operational behavior and system response of the wastewater treatment plant during a specific time period.
\section{Boxplot Analysis}
In this section, we will demonstrate the box-and-line plot analysis of the WWTP operational data. By dividing the data into two parts (Steps 0-244 and Steps 245-609), we generated box line plots for each parameter, and since there are too many parameters, we filtered out the three parameters that have important reference values also known as 'FR', 'SNH', and 'SO'.
\subsection{Flow Rate(FR)}
\begin{figure}[!h]
    \includegraphics[width=5in]{boxplot/FR.png}
    \caption{Boxplot Analysis of FR}
    \label{fig:my_label}
\end{figure}
By observing the box plot of Flow Rate in the above figure, it can be seen that the median of both time ranges is close to the center of the box, indicating that the data for 'FR' is roughly symmetrical, the box is of moderate height, and the whiskers are not too long, suggesting that the data is moderately dispersed. In addition there are no obvious anomalies distributed beyond the ends of the whiskers.
\subsection{Ammonium Nitrogen(SNH)}
\begin{figure}[!h]
    \includegraphics[width=5in]{boxplot/SNH.png}
    \caption{Boxplot Analysis of SNH}
    \label{fig:my_label}
\end{figure}
By observing the box plot of Ammonium Nitrogen in the above figure, it can be seen that the median in the time ranges of 0-244 and 245-609 is skewed to one end of the box, indicating that the distribution of the data is asymmetric and has a certain degree of skewness, and the low height of the box indicates that the data in the FR column has less variation between the 25th and 75th percentiles, which indicates that the data are more stable and has a certain degree of consistency in the data. However, some outliers are observed, which can be further analyzed to see if the cause of these outliers is an error in the data collection process or an abnormality in the operation process.
\subsection{Dissolved Oxygen(SO)}
\begin{figure}[!h]
    \includegraphics[width=5in]{boxplot/SO.png}
    \caption{Boxplot Analysis of SO}
    \label{fig:my_label}
\end{figure}
Through the above observation of the box plot of dissolved oxygen, it can be seen that the median is in the center of the box in the time range of 0-244, indicating that its data are roughly symmetrical, with a high height of the box and a long upper whisker and lower whisker, which usually indicates that the middle 50\% of the data are more loosely distributed, with a high variability among the samples. There are not too many extreme values or outliers, and the overall distribution of the data is relatively even, with no obvious outliers. Whereas the median in the range 245-609 is skewed towards one end of the box, showing some skewness in the distribution of the data, the rest is roughly similar to the first range.
\section{Motif Analysis$^{[6]}$}
In the first part of this section the motif plot detected in the data from the WWTP will be shown and analyzed. The main objective is to identify significant patterns in the data and to analyze the maximum and minimum values observed. This helps us to understand the underlying processes and their changes over time.\\
\\
In this case, reactor 5 in the Folder “close\_loop” is used as an example, and among them, the following parameters can be extracted that have a significant impact on the process: FR, SNH, SNO, SO, TSS, XP.\\
\\
The generated image is divided into two parts, with the upper half showing the raw time series data and the lower half corresponding to the matrix profile. The blue curves in the upper half represent the time series data for each parameter, and the gray rectangles (m1) and light blue rectangles (m2) represent the range of the recognized patterns under two different window lengths (m1 and m2), respectively. The green and orange curves in the lower part represent the window length m1 and the window length m2, respectively, where the green line represents the distance value of the first pattern (m1) in the matrix analysis, and smaller values indicate more similar patterns, and the orange line represents the distance value of the second pattern (m2) in the matrix analysis, and smaller values indicate more similar patterns. The yellow dashed line represents the position of the first pattern, corresponding to the gray rectangular region in the time series data, and the red dashed line represents the position of the second pattern, corresponding to the light blue rectangular region in the time series data.
\subsection{Dissolved Oxygen(SO)}
The parameters in Abbildung 4.4 below represent dissolved oxygen (SO).\\
\begin{figure}[!h]
    \includegraphics[width=5in]{plots/SO.png}
    \caption{Motif Analysis of SO}
    \label{fig:my_label}
\end{figure}
In the upper part of the time series data, the blue curve shows the variation of stream SO concentration. The highest and lowest points in the graph represent the maximum value of 4.395 at index 556 and the minimum value of 0.099 at index 421 in the time series data, respectively.\\
\\
The value of m1 in the graph is 11 hours, which would be a half-day period, while m2 is 24 hours with a one-day period. As can be seen in the lower part of the graph, the maximum values of mp1 and mp2 are 3.0493 (at index 274) and 3.8555 (at index 360), respectively, which implies that at this point in time the pattern of the signal is the least similar to the rest of the time series, and that an anomaly may occur. The minimum value of mp1 is 0.0009 at index 531 and the minimum value of mp2 is 0.0047 at index 412, which means that there is a duplication. In the table, the values of the graphs are shown.\\
\begin{table}[!h]
    \begin{tabular}{c|c|c|c|c}
         & Maximum Value & Minimum Value & Median Value & Standard Deviation \\
         \hline
        SO & 4.39517 & 0.09922 & 0.56416 & 1.2193889
    \end{tabular}
    \caption{Values of SO Parameter}
    \label{tab:my_label}
\end{table}
\subsection{Particulate Organic Matter(XP)}
The following is going to be the parameter XP.\\
\\
\begin{figure}[!h]
    \includegraphics[width=5in]{plots/XP.png}
    \caption{Motif Analysis of XP}
    \label{fig:my_label}
\end{figure}
The first thing that can be seen in the upper part of the time series data is that the blue curve shows the change in particulate organic matter. The highest and lowest points in the graph represent the maximum value of 1209.28 at index 444 and the minimum value of 481.04 at index 313 in the time series data, respectively.\\
\\
In the figure, m1 is 6015h, which gives a period of roughly 250 and a half days, and m2 has a period of roughly 501 days. As can be seen from the dotted lines in the lower part of the figure, the maximum values of mp1 and mp2 are 79.3875 at index 415 and 126.0807 at index 383, respectively, which implies that the signal patterns at these two moments are least similar to the other moments of the time series, and that anomalies may occur. The minimum value of mp1 at index 349 is 39.2495 and the minimum value of mp2 at index 317 is 75.5732, which means that there is a duplication. In the table, the values of the graphs are shown.
\begin{table}[!h]
    \begin{tabular}{c|c|c|c|c}
         & Maximum Value & Minimum Value & Median Value & Standard Deviation \\
         \hline
        XP & 1209.284718 & 481.0434545 & 878.7178501 & 137.609359
    \end{tabular}
    \caption{Values of XP Parameter}
    \label{tab:my_label}
\end{table}
\subsection{Nitrate and Nitrite Nitrogen(SNO)}
The Motif analysis of SNO will be inserted and profiled here.\\
\\
\begin{figure}[!h]
    \includegraphics[width=5in]{plots/SNO.png}
    \caption{Motif Analysis of SNO}
    \label{fig:my_label}
\end{figure}
As can be known from the upper part of the time series data, the blue curve represents the change in SNO. The highest and lowest points in the graph represent the maximum value of 15.29 at index 453 and the minimum value of 2.50 at index 421 in the time series data, respectively.\\
\\
In the figure, m1 is 167h, which gives a period of a week, and m2 has a period of roughly 501 days. As can be seen in the lower part of the graph, the maximum value of mp1 and the maximum value of mp2 are 10.5033 at index 498 and 110.7547 at index 431, respectively, which implies that the signaling pattern at these two moments is least similar to the rest of the time series and may be anomalous. The minimum value of mp1 at index roughly 400 is 1.3554 and the minimum value of mp2 at index 305 is 63.4879, which means that duplication is likely to have occurred in the process. In the table, the values of the graphs are shown.
\begin{table}[!h]
    \begin{tabular}{c|c|c|c|c}
         & Maximum Value & Minimum Value & Median Value & Standard Deviation \\
         \hline
        SNO & 15.289368 & 2.501614 & 7.636299 & 2.1472079
    \end{tabular}
    \caption{Values of SNO Parameter}
    \label{tab:my_label}
\end{table}
\subsection{Ammonium Nitrogen(SNH)}
The next part of the parameter SNH is shown.\\
\\
\begin{figure}[!h]
    \includegraphics[width=5in]{plots/SNH.png}
    \caption{Motif Analysis of SNH}
    \label{fig:my_label}
\end{figure}
In the upper part of the graph, the blue curve represents the change of Ammonium Nitrogen in the wastewater treatment process. The highest and lowest points in the graph represent the maximum value of 12.59 at index 256 and the minimum value of 0.04 at index 448 in the time series data, respectively.\\
\\
The value of m1 in the graph is 11 hours, which would be a half-day period, while m2 is 24 hours with a one-day period. From the lower part of the graph, information can be derived that the maximum value of mp1 and the maximum value of mp2 are 2.9869 at index 523 and 2.6270 at index 523, respectively, which implies that the signaling pattern at these two moments is least similar to the rest of the time series and may be anomalous. The minimum value of mp1 at index 267 is 0.0006 and the minimum value of mp2 at index 514 is 0.0060, which means that a duplication has occurred. In the table, the values of the graphs are shown.
\begin{table}[!h]
    \begin{tabular}{c|c|c|c|c}
         & Maximum Value & Minimum Value & Median Value & Standard Deviation \\
         \hline
        SNH & 12.591682 & 0.046862 & 0.5103675 & 1.8107844
    \end{tabular}
    \caption{Values of SNH Parameter}
    \label{tab:my_label}
\end{table}
\subsection{Total Suspended Solids(TSS)}
The following is a Motif analysis of the parameter TSS, also known as Total Suspended Solids.\\
\\
\begin{figure}[!h]
    \includegraphics[width=5in]{plots/TSS.png}
    \caption{Motif Analysis of TSS}
    \label{fig:my_label}
\end{figure}
In the upper part of the graph, the blue curve can be seen, which represents the variation of Total Suspended Solids during the wastewater treatment process. The highest and lowest points in the graph represent the maximum value of 4556.32 at index 584 and the minimum value of 2540.66 at index 256 in the time series data, respectively.\\
\\
In the figure, m1 is 6015h, which gives a period of roughly 250 and a half days, and m2 has a period of roughly 501 days. From the lower part of the graph, it can be seen that the maximum value of mp1 and the maximum value of mp2 are 88.8894 at index 432 and 134.0381 at index 409, respectively, which implies that the signaling patterns of these two moments are least similar to the other moments of the time series, and that anomalies may exist. The minimum value of mp1 at index 463 is 52.7555 and that of mp2 at index 353 is 84.3578, which means that there is a high degree of repetition. In the table, the values of the graphs are shown.
\begin{table}[!h]
    \begin{tabular}{c|c|c|c|c}
         & Maximum Value & Minimum Value & Median Value & Standard Deviation \\
         \hline
        TSS & 4556.323702 & 2540.655701 & 3785.580318 & 352.7483605
    \end{tabular}
    \caption{Values of TSS Parameter}
    \label{tab:my_label}
\end{table}
\subsection{Flow Rate(FR)}
The following figure shows the results of the motif analysis in the time series data for FR.\\
\\
\begin{figure}[!h]
    \includegraphics[width=5in]{plots/FR.png}
    \caption{Motif Analysis of FR}
    \label{fig:my_label}
\end{figure}
It can be seen that Motif 1 and Motif 2 correspond to the lowest points in the matrix profile, respectively, which implies that the subsequences at these positions have the highest similarity in the whole time series. In the upper part of the time series data, the blue curve shows the variation of Flow Rate. The highest and lowest points in the graph represent the maximum value of 142872 at index of 256 and the minimum value of 88038 at index of 452 for the time series data, respectively.\\
\\
The value of m1 in the graph is 11 hours, which would be a half-day period, while m2 is 24 hours with a one-day period. In the lower part of the figure, the maximum value of mp1 and the maximum value of mp2 are 2.5018 at index 313 and 3.7598 at index 529, respectively, which implies that at this point in time, the pattern of the signal has the lowest similarity with the rest of the time series, and an anomaly may occur. At index 351, the minimum value of mp1 is 0.0001, and at index 452, the minimum value of mp2 is 0.0393, which implies that at this point in time, the pattern of the signal has the highest similarity to the rest of the time series, and may indicate the occurrence of a common or repeating event. In the table, the values of the graphs are shown.
\begin{table}[!h]
    \begin{tabular}{c|c|c|c|c}
         & Maximum Value & Minimum Value & Median Value & Standard Deviation \\
         \hline
        FR & 142872.6052 & 88038.55187 & 102685.1505 & 8386.555512
    \end{tabular}
    \caption{Values of FR Parameter}
    \label{tab:my_label}
\end{table}
\section{Motif Analysis of Periodic Patterns Over 7 Days}
In order to have a clearer as well as accurate view of cyclic behavior and anomalies in the data, the index was shortened to 7, so this section is for the theme extraction process for the 7-day time series data from the WWTP. The objective is to identify recurring patterns (themes) in the three parameters that are more important, i.e., FR (flow rate), SNH (soluble ammonium), and SO (soluble oxygen).\\
\\
The generated plots show the time series data as well as the matrix profiles for each parameter. \\
The gray rectangles represent the patterns determined from the calculation of the period m1, where the period m2 is omitted as it is comparable to the period m2.\\
The yellow dashed line represents the time index of the minimum value in the matrix profile, indicating the most frequently occurring pattern. The red dashed line represents the time index of the maximum value in the matrix profile, highlighting significant deviations or anomalies.
\subsection{Motif and Matrix Profile for FR (Flow Rate)}
As seen in the figure the value of m1 is 12h, which would represent a period of half a day. The recurring pattern that can be identified through the image is the gray rectangle at the top of the figure, which marks the pattern (theme) that occurs most frequently in the stream data. This pattern represents typical behavior that occurs periodically over a 7-day period. Understanding this pattern will help optimize flow control and ensure stable and efficient operations.\\
\\
In contrast, anomalies appear in the red dotted line at the bottom of the chart, which represents the time index with the highest median in the matrix profile, indicating an anomaly or deviation from the typical pattern. This requires further investigation to determine the cause of the irregularity, which may be due to equipment failure or unexpected changes in inflow. It is also known from the figure that the maximum value is in the matrix profile for FR with m1=12 is located at day 5.7917. Its maximum and minimum values will be shown in the table.
\begin{figure}[!h]
    \includegraphics[width=5in]{plots2/FR.png}
    \caption{Motif Analysis in FR over 7 days}
    \label{fig:my_label}
\end{figure}
\begin{table}[!h]
    \begin{tabular}{l|c|r}
         & Maximum Value & Minimum Value \\
         \hline
        FR & 117976.3493  & 90492.95868
    \end{tabular}
    \caption{FR Parameter Statistical Indicator}
    \label{tab:my_label}
\end{table}
\subsection{Motif and Matrix Profile for SNH (Soluble Ammonium)}
\begin{figure}[!h]
    \includegraphics[width=5in]{plots2/SNH.png}
    \caption{Motif Analysis in SNH over 7 days}
    \label{fig:my_label}
\end{figure}
The regularity of ammonium levels can be derived from the graph, and the themes identified for SNH show typical fluctuations in ammonia levels within the plant. By recognizing this pattern, the dosage of chemicals used for ammonia removal can be better managed, thus improving overall treatment efficiency.\\
\\
As seen in the figure the value of m1 is 12h, which would represent a period of half a day. Significant deviations are present in the yellow dashed line at the bottom of the graph, indicating time periods when ammonia levels deviate from normal. This deviation speculation may indicate a problem with the biological treatment process or a change in influent characteristics that requires corrective action. It can be also learned from the figure that the maximum value is in the matrix profile for SNH with m1=12 is located at day 6.2292. Its maximum and minimum values will be shown in the table.
\begin{table}[!h]
    \begin{tabular}{l|c|r}
         & Maximum Value & Minimum Value \\
         \hline
        SNH & 5.678577818  & 0.061569123
    \end{tabular}
    \caption{SNH Parameter Statistical Indicator}
    \label{tab:my_label}
\end{table}
\subsection{Motif and Matrix Profile for SO (Soluble Oxygen)}
\begin{figure}[!h]
    \includegraphics[width=5in]{plots2/SO.png}
    \caption{Motif Analysis in SO over 7 days}
    \label{fig:my_label}
\end{figure}
SO in the graph reflects the normal oxygen supply required to maintain optimal biological activity in the treatment plant, and consistently meeting this pattern ensures that aerobic bacteria are effective in decomposing organic matter. And as seen in the figure the value of m1 is 12h, which would represent a period of half a day.\\
\\
Maximum values in the matrix profile (marked with a red dashed line) reflect irregularities in the oxygen supply, which is an indication that oxygen levels are not being maintained as expected. So the maximum value in the graph is in the matrix profile for SO with m1=12 is located at day 0.4688, which is about 11 hours out. Its maximum and minimum values will be shown in the table.
\begin{table}[!h]
    \begin{tabular}{l|c|r}
         & Maximum Value & Minimum Value \\
         \hline
        SO & 3.524592066  & 0.152796907
    \end{tabular}
    \caption{SO Parameter Statistical Indicator}
    \label{tab:my_label}
\end{table}
\section{Motif Analysis of Periodic Patterns Over 250 Days}
This section focuses on three important parameters in the reactor dataset: SNO, TSS, and XP. The cyclical characteristics of these three parameters are examined in depth to identify patterns and trends over a 250-day period.\\
\\
The SNO, TSS and XP columns were first extracted from the dataset and the time span was set to 250 days. As above, there are 4 data points per hour, so 96 data points per day, and the time series plots and matrix contours for each of the three parameters will be analyzed in detail next.
\subsection{Motif and Matrix Profile for SNO (Nitrate and Nitrite Nitrogen)}
\begin{figure}[!h]
    \includegraphics[width=5in]{plots2/SNO.png}
    \caption{Motif Analysis in SNO over 250 days}
    \label{fig:my_label}
\end{figure}
As shown, the value of m1 is 166 hours, so the period is roughly one week. Maximum values show sections with less similarity and may reflect anomalies or periods of high variability. The range of maximum values is shown below with the help of a table.\\
\begin{table}[!h]
    \begin{tabular}{l|c|r}
         & Maximum Value & Max Matrix Profile Index(days) \\
         \hline
        SNO & 14.299288  & 58.791667
    \end{tabular}
    \caption{SNO Parameter Statistical Indicator}
    \label{tab:my_label}
\end{table}
From the table it can be seen that at the maximum values then the less similar parts are shown, possibly reflecting anomalies or periods of high variability.
\subsection{Motif and Matrix Profile for Total Suspended Solids(TSS)}
\begin{figure}[!h]
    \includegraphics[width=5in]{plots2/TSS.png}
    \caption{Motif Analysis in TSS over 250 days}
    \label{fig:my_label}
\end{figure}
As seen in the figure the value of m1 is 12h, which would represent a period of half a day. An analysis similar to SNO performs the same processing and visualization of the TSS data. By looking at the time series plots and matrix contour plots of the TSS, periodic features and anomalies can be identified.
\begin{table}[!h]
    \begin{tabular}{l|c|r}
         & Maximum Value & Max Matrix Profile Index(days) \\
         \hline
        TSS & 4530.919377  & 197.843750
    \end{tabular}
    \caption{TSS Parameter Statistical Indicator}
    \label{tab:my_label}
\end{table}
\subsection{Motif and Matrix Profile for Organic Matter(XP)}
\begin{figure}[!h]
    \includegraphics[width=5in]{plots2/XP.png}
    \caption{Motif Analysis in XP over 250 days}
    \label{fig:my_label}
\end{figure}
As seen in the figure the value of m1 is 3000h, which would represent a period of 125 days. By looking at the time series plot and the matrix contour plot of XP above, its periodic characteristics and outliers can be similarly identified and the maximum value obtained.
\begin{table}[!h]
    \begin{tabular}{l|c|r}
         & Maximum Value & Max Matrix Profile Index(days) \\
         \hline
        XP & 1203.552444  & 30.645833
    \end{tabular}
    \caption{XP Parameter Statistical Indicator}
    \label{tab:my_label}
\end{table}
\chapter{Conclusion}
In this study focuses on the detailed analysis of the key parameters of the wastewater treatment plant through Fourier transform and matrix profile analysis. Taking an example of dissolved oxygen (SO) parameter from the 7-day Motif analysis, in the analyzed results, we found that its maximum frequency corresponds to a period of about 24 hours, which indicates that the SO parameter shows significant 24-hour cyclic fluctuations. This cyclic fluctuation suggests that there are regular changes in DO levels at specific times of the day. This finding is of great practical significance as it suggests that aeration needs to be adjusted at specific times of the day to maintain proper oxygen levels and ensure stable system operation.\\
\\
Matrix profile analysis, on the other hand, helps us identify significant patterns (motifs) in the time series that recur prior to an abnormal run and can be used as early warning signals. For example, the ammonia nitrogen (SNH) parameter, which has a significant peak at a certain point in time, represents a specific pattern that occurs prior to an abnormal emission event, which can be used to take early action.\\
\\
By analyzing the maximum and minimum values of different parameters in the above section, it can be observed that these values reflect the operational status of the system during periods of high and low loads. For example, the maximum value of flow rate (FR) reveals a peak in water intake, whereas the minimum value found in the graph of SO may indicate a situation of insufficient oxygen supply, which may affect the activity of microorganisms.\\
\\
This study reveals the periodicity and significant patterns of key parameters in wastewater treatment plants, aiming to provide valuable insights into the optimal operation of the system, which can be continually improved through further research and data analysis to improve the treatment efficiency and stability of wastewater treatment plants.
\chapter{Appendlces}
\begin{figure}[!h]
    \includegraphics[width=5in]{Appendlces/heatmap.png}
    \caption{Heatmap of Parameters$^{[7]}$}
    \label{fig:my_label}
\end{figure}
\begin{figure}[!h]
    \includegraphics[width=5in]{Appendlces/FR.png}
    \caption{Discord Analysis$^{[8]}$ in FR}
    \label{fig:my_label}
\end{figure}
\begin{figure}[!h]
    \includegraphics[width=5in]{Appendlces/TSS.png}
    \caption{Discord Analysis in TSS}
    \label{fig:my_label}
\end{figure}
\begin{figure}[!h]
    \includegraphics[width=5in]{Appendlces/XP.png}
    \caption{Discord Analysis in XP}
    \label{fig:my_label}
\end{figure}
\chapter{References}
$[1]$ https://www.britannica.com/technology/wastewater-treatment \\
$[2]$ https://pandas.pydata.org/pandas-docs/stable/reference/api/pandas.DataFrame.html \\
$[3]$ https://stumpy.readthedocs.io/en/latest/Tutorial\_STUMPY\_Basics.html \\
$[4]$ https://seaborn.pydata.org/generated/seaborn.boxplot.html \\
$[5]$ https://matplotlib.org/stable/gallery/index.html \\
$[6]$ https://stumpy.readthedocs.io/en/stable/Tutorial\_STUMPY\_Basics.html\#find-a-motif-using-stump \\
$[7]$ https://seaborn.pydata.org/generated/seaborn.heatmap.html \\
$[8]$ https://stumpy.readthedocs.io/en/stable/Tutorial\_STUMPY\_Basics.html\#find-potential-anomalies-discords-using-stump
\end{document}
